\section{A residual tensor presentation on the prescribed Fontaine comparison}

% BEGIN R80D1
\paragraph{R80D1. The reduced coefficient ring and its letters.}
Fix an original $\mathrm{TopSit}$, one of its residue fields
$k$ with $O\twoheadrightarrow k$ and kernel $(\lambda)$, and
the given local Wach and Fontaine data at layer one. Put
$\mathcal E=W.\mathrm{EE}$, $U=O\otimes_{\mathbb Z_\ell}
W(C^\flat)$ and $B=k\otimes_O U$. Here $C^\flat$ denotes the
fraction field of the integral tilt. The tie gives
$\mathcal E\longrightarrow U$, and we use its actual image on
reducing to $B$. Let $e_i$ and $\epsilon_i$ denote the images
in $\mathcal E$ and $B$ of the indicator of the $i$th letter.
They are orthogonal idempotents summing to $1$.

Each $\epsilon_i$ is nonzero. The inclusion of the letter ring
in $\mathcal E$ is injective by the Wach datum. The comparison
embedding of $\mathcal E$ in $U$ is injective: a scalar in its
kernel annihilates the injectively compared free module of rank
one attached to the trivial lattice. An idempotent in $\lambda U$
lies in $\lambda^jU$ for every $j$ by taking powers. This
intersection is zero since $\ell$ is a unit times a positive
power of $\lambda$ and $U$ is finite free over the separated
Witt ring (use a $\mathbb Z_\ell$-basis of $O$). Finally
\[
 B=k\otimes_{\mathbb F_\ell} C^\flat,\qquad B^{\varphi=1}=k .
\]
Indeed reduction of the Witt vectors of the perfect tilt modulo
$\ell$ is their zeroth Witt coordinate. Frobenius on this tensor
is the identity on $k$ and $x\mapsto x^\ell$ on $C^\flat$.
The kernel of $x\mapsto x^\ell-x$ on that field is
$\mathbb F_\ell$, and tensoring its kernel with the vector space
$k$ preserves exactness. The Galois action comes from the actual
action on the tilt, commutes with Frobenius and fixes $k$.

Write $F=k((X))$. There is a surjection
$\mathcal E\to F$ at each letter, sending $e_i a$ to the
reduction of $a_i$ for a tuple of power series $a$. More precisely
$e_i\mathcal E/\lambda e_i\mathcal E=F$. To see it, clear a
power of the invertible Wach variable $X$ in an arbitrary element
of $\mathcal E$ modulo $\ell$ using the density field of $W$.
The saturation field says that a tuple of power series divisible
by $\ell$ in $\mathcal E$ is already divisible there in the
power series ring. Thus evaluation of the resulting Laurent
series modulo $\lambda$ is well-defined, using common powers to
compare two representatives. Conversely a lift coefficient by
coefficient, followed by division by the same power of $X$, gives
every Laurent series. When its $i$th reduction is zero the
$i$th numerator is coefficientwise divisible by $\lambda$;
the error is divisible by $\ell$, hence by $\lambda$ as well.
This proves the kernel assertion. It gives field embeddings
$F\to\epsilon_i B$ since these component rings are nonzero.
Their sum is the diagonal homomorphism $\kappa:F\to B$.
Frobenius carries the preceding letter to the current one and
sends $X$ to $X^\ell$ on the reductions. A Galois lift fixes
each idempotent and acts by the reduction of its cyclotomic
substitution. These identities hold on all Laurent series by
lifting power series and inverting $X$.
% END R80D1

% BEGIN R80D2
\paragraph{R80D2. The fixed Laurent parameter, Frobenius order and generator.}
Here $K$ is unramified and $\ell\ge5$. The cyclotomic action
on the tower quotient $\Gamma$ identifies it with
$\mathbb Z_\ell^\times$: the kernel has already been divided
out, and the totally ramified cyclotomic tower is disjoint from
$K$. Thus the finite torsion group $\Delta$ is the Teichm\"uller
group of order $\ell-1$, and we can choose $\gamma$ with
$\chi_\Gamma(\gamma)=1+\ell$. The reductions of their actions
are respectively $X\mapsto(1+X)^a-1$ for $a\in\Delta$ and
$X\mapsto(1+X)^{1+\ell}-1$. The binomial series identities
follow by continuity degree by degree from integral exponents.
Write
\[
        v=1+\sum_{a\in\Delta}(1+X)^a\quad\text{in }k[[X]].
\]
This is the reduction of the native $W.\pi_{01}$, since its
additional multiple of $\ell$ vanishes. Torsion substitutions
permute the summands. In degrees $<\ell$ the exponents have all
nonzero residues, so the initial segment of $v$ agrees with
$\sum_{j=0}^{\ell-1}(1+X)^j=X^{\ell-1}$. For an invariant
Laurent series its leading degree is divisible by $\ell-1$,
by applying an element whose linear coefficient generates
$\mathbb F_\ell^\times$. W80N1 identifies the integral fixed ring in this parameter.
Multiplication by a power of $v$ clears a Laurent pole. Hence
$F^\Delta=k((v))$ and $k[[X]]^\Delta=k[[v]]$.
The linear coefficients distinguish all the substitutions.

Take the residual scalar field $F_0=k((Z))$ and
$\iota(Z)=v$. Frobenius gives $v(X^\ell)=v(X)^\ell$, as the
coefficients of $v$ lie in the prime field. Put
$h(X)=X+X^\ell+X^{\ell+1}$, the substitution for $\gamma$.
From $v=X^{\ell-1}+O(X^\ell)$ one obtains
$v(h)-v=-X^{2(\ell-1)}+O(X^{2\ell-1})$:
in the leading monomial choose one occurrence of $h-X$, whose
coefficient is $\ell-1=-1$; higher monomials and all terms with
two such occurrences have higher order. This difference is
torsion invariant, hence
$v(h)=v(1-b(v)v)$ with $b\in k[[Z]]$, $b(0)=1$.
Transport Frobenius and $\gamma$ to $F_0$ along $\iota$.
They commute, preserve integral series, act on $Z$ by $Z^\ell$
and $Z(1-b(Z)Z)$, and the second is invertible by substitution
of a power series with linear coefficient $1$. These are the
original $\mathrm{RingPG}$ and cyclotomic conditions, including
the precise normalization of $\iota(Z)$.

Let $\tau(i)$ be the unique Wach label whose embedding is the
given embedding at $i$. The two embedding lists are bijections
and satisfy the same predecessor formula by precomposition with
the unramified Frobenius. Thus $\tau$ is bijective and
$\tau(i-1)=\tau(i)-1$. The lifts of $\Delta$ and of $\gamma$
generate a dense subgroup of $G_K$: they include the kernel of
$G_K\to\Gamma$ and their images generate densely, since
$1+\ell$ topologically generates $1+\ell\mathbb Z_\ell$.
A quotient of an original lattice by a submodule containing
$\lambda T$ has continuous action on its finite discrete space.
Consequently a vector or map between such representations fixed
by these lifts is fixed by the whole group.
% END R80D2

% BEGIN R80D3
\paragraph{R80D3. Comparison of a lattice quotient and its exact integral kernel.}
Let $q:T\twoheadrightarrow A$ be an equivariant quotient of an
original lattice with $\dim_k A=m$. For a Wach letter $i$ set
$D_i(T)=e_iD(T)/\lambda e_iD(T)$, an $F$-space of dimension
$d=\operatorname{rank}_O T$. Indeed the source's freeness and
positive finite rank give a finite basis when $d>0$; when
$d=0$ the comparison embeds $D(T)$ into zero. On the reduced
comparison use
$J_q=(1\otimes q)\overline{j_T}$ and let $V_i(q)$ be the
image of $D_i(T)$ under $\epsilon_i J_q$ in
$B_i\otimes_k A$, where $B_i=\epsilon_i B$. In bases of the two
free modules the original comparison over $U$ is a square matrix
whose columns span $U\otimes_O T$, by the tie. A surjective square
matrix over a commutative ring is invertible, by a right inverse
and determinants. Thus scalar extension gives an isomorphism
\[
       B_i\otimes_F D_i(T)\ \simeq\ B_i\otimes_O T .
\]
It carries every map of lattices to its ordinary base change,
by naturality of the tie.

Take $T'$ with an injective lattice morphism $f:T'\to T$
having image $\ker q$. Such a lattice always exists by using the
kernel inside the same rational space, since $\lambda T\subseteq
\ker q$. For any prescribed such $f$, the image of its reduction
$T'/\lambda T'\to T/\lambda T$ has dimension $d-m$. By the square
comparison, $D_i(f)$ has this same rank over $F$: changes of basis
over $B_i$ preserve determinantal ideals, and a rank-$s$ matrix
over either field has unit determinantal ideals through size $s$
and zero ones thereafter in the nonzero ring $B_i$. The map
from $D_i(T)$ to $B_i\otimes_k A$ kills this image. Its image
after $B_i$-span is all of $B_i\otimes_k A$, so it has dimension
at least $m$ over $F$ (reduce any spanning family modulo a maximal
ideal of $B_i$). Rank-nullity now gives dimension exactly $m$
and kernel precisely $\operatorname{im}D_i(f)$.
In particular an $F$-basis of $V_i(q)$ is also a $B_i$-basis of
$B_i\otimes_k A$, again by the square matrix argument.

The integral kernel formula requires an equality before reduction.
For $x\in D(T)$ mapping to zero on this letter, lift its residual
preimage to $y\in e_iD(T')$. Then
$e_i x-D(f)y=\lambda z$ in $e_iD(T)$. The map $s=f^{-1}\lambda$
is a lattice morphism $T\to T'$ since $\lambda T\subset f(T')$;
it is defined rationally on the whole space since $f(T')$ spans it.
Functoriality and the scalar law give $D(f)D(s)z=\lambda z$.
Thus $e_ix=D(f)(y+D(s)z)$. Conversely any such equality kills
the letter image by $qf=0$. This is precisely the unaveraged
kernel witness required by the native presentation, also for
$x$ that happens to be $\Delta$-fixed.
% END R80D3

% BEGIN R80D4
\paragraph{R80D4. Descent through torsion and independence of the presenting lattice.}
A lift of $\delta\in\Delta$ acts on $V_i(q)$ by acting on the
reduced comparison tensor. On this subspace that action agrees
with $\gamma_D(\delta)$ on representatives, so it is independent
of the lift. The action is semilinear for the faithful torsion
action on $F$ in R80D2. Its fixed subspace $V_i(q)^\Delta$
has an $F_0$-basis that is an $F$-basis of $V_i(q)$.
Indeed sums $\sum_\delta\delta(cv)$, for $c\in F$ and
$v\in V_i(q)$, are fixed and span $V_i(q)$ over $F$.
Otherwise a nonzero $F$-linear functional annihilating their
span would give $\sum_\delta \delta(c)u(\delta v)=0$ for every
$c$; linear independence of the distinct field automorphisms
forces $u(v)=0$ for all $v$. Choose a basis from the fixed vectors.
Uniqueness of coefficients then identifies its fixed coefficients
with $\iota(F_0)$. Every fixed vector lifts to $D(T)^\Delta$:
average a representative in $e_iD(T)$ over $\Delta$, multiplying
by a lift to $O$ of $|\Delta|^{-1}\in k$. This does not change
its residual image.

For a $G_K$-map $a:A\to A'$ between two presented objects
$q:T\twoheadrightarrow A$, $q':T'\twoheadrightarrow A'$, use
the lattice of pairs $(u,u')$ with $a q(u)=q'(u')$. It is full in
$T\oplus T'$, contains $\lambda(T\oplus T')$, is stable and its
first projection is onto $T$. Applying the square comparison of
R80D3 to that projection shows that the induced map on the
$F$-spaces $D_i$ is onto (the ordinary reduction is onto, so
the maximal minors generate the unit ideal). Lifting in that
space and applying the second projection shows that $1\otimes a$
carries $V_i(q)$ into $V_i(q')$. The same holds on invariants.
For $a=1_A$ use both directions: the two subspaces, hence their
fixed subspaces, are equal inside $B_i\otimes_k A$ for every
presentation of $A$. Write $V_i(A)$ for this common space.
The allowed objects are closed under subobjects (take the preimage
lattice), quotients (compose the surjection) and jointly surjective
sums (use the sum of the two presenting lattices). They include
the trivial line via reduction of the unit lattice.
% END R80D4

% BEGIN R80D5
\paragraph{R80D5. The exact functor and full faithfulness on the presented objects.}
Use $V_{\tau(i)}(A)^\Delta$ with its $F_0$-basis from R80D4 as
the $i$th copy of $F_0^m$. Under this embedding into
$B_{\tau(i)}\otimes_k A$, define Frobenius by acting on the first
tensor factor, the chosen $\gamma$ by any lift acting on both
factors, and $F(a)_i$ by $1\otimes a$. These maps preserve the
specified spaces by the tie's intertwining identities and R80D4.
The operators are semilinear for R80D2, Frobenius shifts letters
by $+1$, and the operators commute. The action of $\gamma$ has
its inverse from a lift of $\gamma^{-1}$. The Frobenius images of
a basis span on the next letter: that basis is an invertible square
matrix over $B_{\tau(i-1)}$, so its image under the unital letter
Frobenius is invertible over $B_{\tau(i)}$. Expressing these columns
in the chosen target frame makes the determinant a nonzero element
of $F_0$. In the same way the rank of $F(a)_i$ over $F_0$ equals
the rank of $a$ over $k$ by comparison over the nonzero ring
$B_{\tau(i)}$. Identities, composition, addition and scalar
multiplication follow inside the comparison tensors. This gives
the native exactness on injective and surjective maps. For objects
outside the allowed class use the coordinatewise scalar Frobenius
and $\gamma$, and use the zero map where a map law has no
hypothesis. The unit at $i$ is the frame coordinate of
$\epsilon_{\tau(i)}\otimes1$; it is nonzero, fixed under $\gamma$
and carried to the next unit by Frobenius.

In fact every family of $F_0$-linear maps between these presented
objects commuting with both operators comes from a unique $G_K$-map.
Extend it in the two compared bases on each letter to a
$B_{\tau(i)}$-linear matrix, and combine the matrices by the
orthogonal idempotents to a matrix $M$ over $B$. Frobenius
compatibility implies $\varphi(M)=M$: check on the Frobenius
images of a frame, which form a basis in the next comparison tensor.
R80D1 gives $M$ over $k$. All lifts of $\Delta$ fix the frame
vectors on both sides; this is the equality on representatives
in R80D4. The same comparison with the $\gamma$ action verifies
the intertwining identity for a lift of $\gamma$. Since the
entries are constants and the frames span, $M$ intertwines the
original representations on these lifts, hence everywhere by
R80D2. Its image under the functor has the given components,
and uniqueness follows by the same scalar extension. In particular
a compatible fixed vector at each letter is the image of the
unit family along a morphism from the trivial representation.
This proves the fixed-vector clause with its specified unit.
% END R80D5

% BEGIN R80D6
\paragraph{R80D6. The presentation maps, the unit lattice and the return to the given situation.}
At a surjection $q:T\twoheadrightarrow A$, define $j_{q,i}(x)$
for $x\in D(T)^\Delta$ as the coordinates of
$\epsilon_{\tau(i)}J_q(x)$ in the R80D4 invariant frame.
Define it to be zero for other $q$. It is additive and onto in
the required case by the averaging in R80D4. If a letter scalar
$a$ has $\overline{a_{\tau(i)}}=\iota(c)$ for $c\in k[[Z]]$,
multiplication by $a$ on representatives induces multiplication
by $c$ on coordinates whenever both representatives are invariant.
R80D5 and naturality of $J_q$ give exactly its Frobenius, $\gamma$
and map identities. R80D3 gives its integral kernel formula with
the idempotent at $\tau(i)$.

For the unit, take the tie's generator $u$ of its Wach lattice
mapping to $1\otimes u_0$, with $u_0$ a basis of the trivial
$O$-line. It is $\Delta$-invariant by injectivity of comparison.
Every other $\Delta$-invariant generator is $a u$, with $a$ a
unit of the letter power series ring: the converse multiple gives
an inverse since a scalar annihilating $u$ would annihilate the
nonzero free generator after the injective coefficient embedding.
Invariance gives $\delta(a)=a$ by the same cancellation.
Each reduction of $a$ and of $a^{-1}$ therefore lies in $k[[v]]$.
The image of $u$ at a surjection to the trivial line is the unit
of R80D5 times the nonzero scalar $q(u_0)$. Thus $j_{q,i}(au)$
and the specified unit generate exactly the same $k[[Z]]$-lattice.
Together with $\iota(Z)=\overline{W.\pi_{01}}$ these are all the
fields of $\mathrm{Pres}$. R80D2 supplies the required matched
embedding at every letter and $\chi_\Gamma(\gamma)=1+\ell$.
This proves $S.\mathrm{RealD}(H)$ for every native $S$ and every
one of its residue fields on the given Fontaine comparison.
L78W13 supplies that comparison for the project's chosen data
under its original integral crystalline input.

The same construction gives the abstract tensor-presentation
implication of the source. In that formulation take
$B=k\otimes_O U$ for its given comparison ring $U$. Its hypotheses
give the scalar Frobenius fixed ring, detection of idempotents,
the commuting group action, and the cyclic data directly.
Use its $\ell$-adic density and saturation of $\mathcal E$ as in
R80D1, the supplied faithful Teichm\"uller actions and cyclotomic
fixed parameter as in R80D2, and R80D3--R80D4 unchanged.
For the required fixed-vector clause the abstract density premise
is only about vectors of presented objects. Sum the components of
a compatible vector using the idempotents: Frobenius fixes the
sum, so it is $1\otimes t$. The component frames show that $t$
is fixed by every torsion lift and by every lift of $\gamma$.
The abstract premise now gives $t\in A^G$. Sending $1$ to $t$
gives the required morphism from the unit, without any additional
density premise on general Hom spaces. All other presentation
identities follow by the same matrix comparisons.
% END R80D6
